\section{Contribution to the Dissertation}
\label{sec:sources-wasm4pm-contribution}

\subsection{Contribution to Prediction vs.~Coordination}
\label{sec:sources-wasm4pm-contrib-prediction-coordination}

The dissertation's central organizing distinction is between prediction and coordination:
prediction answers ``what will happen'' while coordination answers ``what must happen and did it.''
Process-mining tools in the old regime are primarily predictive instruments — they discover
models and compute fitness scores, but those scores are advisory rather than enforceable.
\texttt{wasm4pm} shifts the register from prediction to coordination. The 0.8 admission threshold
in the \texttt{ConformanceVerdict::admits()} method is not a recommendation; it is a gate. A
trace that does not meet the threshold is refused with a typed \texttt{ConformanceRefusal} and
cannot advance to the next lifecycle state. The six Blue River Dam proof gates in the Lifecycle
Authority make this coordination role explicit: advancement from Design to Simulation requires a
passing \texttt{SoundnessProof}; advancement from Simulation to MonitoringOps requires a passing
\texttt{ConformanceProof}. The typestate machine enforces these conditions at compile time where
possible and at runtime where the proof must be computed. This is the dissertation's primary
claim instantiated in executable Rust: coordination requires proof, and proof requires receipts.
[AUTHOR THESIS]

\subsection{Contribution to Process Evidence}
\label{sec:sources-wasm4pm-contrib-process-evidence}

\texttt{wasm4pm} makes two concrete contributions to the process-evidence strand of the
dissertation. First, it demonstrates that the \texttt{Evidence<T,~State,~W>} lattice can be
implemented in a zero-external-dependency Rust crate with compile-time enforcement of
monotonicity: witness markers implement \texttt{Lattice<T>} with provably monotone \texttt{join}
operations, and admission state transitions are enforced by the \texttt{Admitted} enum (no
transition from \texttt{Sealed} back to \texttt{Initial} is possible at the type level). Second,
it demonstrates that the four mining algorithms (Alpha, Inductive, Heuristics, DFG) can each
emit typed, lattice-compliant witnesses rather than untyped floating-point outputs. The seven
inline unit tests in \texttt{src/mining/mod.rs} — verified by the sealed
\texttt{MINING\_RENDER\_RECEIPT.md} — constitute the first receipted proof that this evidence
lattice is both constructible and testable in the process-intelligence context. The A* alignment
engine's use of a 5000-iteration cap and \texttt{BinaryHeap} demonstrates that optimal alignment
is achievable within a bounded computational budget, addressing the scalability concern raised in
the conformance-checking literature.

\subsection{Contribution to Manufacturing Architecture}
\label{sec:sources-wasm4pm-contrib-manufacturing}

The mining module's provenance demonstrates the dissertation's manufacturing architecture
end-to-end. The module was not written by hand; it was materialized by the ggen manufacturing
machinery from templates bound to authoritative specifications. The sealed receipt
(\texttt{MINING\_RENDER\_RECEIPT.md}) documents the complete manufacturing chain, and the
SHA-256 content hash \texttt{08a067d1ee19ea67150c194e9a2db7d86dfd994223b922de7e8606f45fbdf8e5}
binds the artifact to that chain. This demonstrates that the manufacturing architecture is not
merely a documentation claim: it produces verifiable, compilable, tested Rust source. The
\texttt{verify\_jcs\_receipt\_signature} function in \texttt{src/crypto.rs} (JCS RFC 8785
canonicalization + Ed25519 verification) provides the cryptographic primitive for verifying that
receipt JSON envelopes were not tampered with after issuance. Together, these contributions
establish that the manufacturing architecture can produce receipted artifacts whose provenance is
auditable without trusting any external party.

\subsection{Contribution to Capital-Flow Infrastructure}
\label{sec:sources-wasm4pm-contrib-capital-flow}

The dissertation's capital-flow strand — connecting process evidence to settlement, DAO
governance, and capital allocation decisions — requires that process fitness scores be
tamper-evident and attributable. \texttt{wasm4pm}'s Conformance Authority contribution is that
every fitness score, once signed and chained, cannot be retroactively modified without invalidating
the receipt chain. The spot-audit framework (specified in \texttt{research-verdict.md}) provides
probabilistic tamper detection: a verifier can re-replay any archived trace against the model
snapshot and confirm the receipt. In a capital-flow context, this means that a claim of
``process fitness $\geq 0.95$'' can be backed by a verifiable receipt chain rather than an
assertion. The \texttt{DecommissioningReceipt} type and the \texttt{execute\_oblivion\_protocol}
function contribute the archival finality layer: once a process model is decommissioned, the
receipt chain is closed and the final archival proof links backward through the complete history.
[INTERPRETATION] This positions \texttt{wasm4pm} as a foundational primitive for
process-evidence-backed capital allocation decisions, where settlement conditions can be
conditioned on verifiable process conformance rather than self-reported metrics.
